% ECONOMICS_PROBLEM: {"id": "J23-400", "title": "AI and automation: displacement versus complementarity", "jel_code": "J23", "source_id": "OP-0022"}
% FIRST_POSTED: 2026-10-09
% ATTEMPT_STATUS: unresolved
% ENTRY_KIND: research_attempt
\documentclass[11pt]{article}
\usepackage[T1]{fontenc}
\usepackage[utf8]{inputenc}
\usepackage[margin=25mm]{geometry}
\usepackage{amsmath,amssymb,amsthm,xurl,hyperref}
\newtheorem{proposition}{Proposition}
\newtheorem{lemma}{Lemma}
\newtheorem{corollary}{Corollary}
\newcommand{\E}{\mathbb E}
\newcommand{\R}{\mathbb R}
\newcommand{\Var}{\operatorname{Var}}
\newcommand{\Cov}{\operatorname{Cov}}
\DeclareMathOperator*{\argmax}{arg\,max}
\setlength{\parindent}{0pt}
\setlength{\parskip}{6pt}
\title{AI and automation: displacement versus complementarity: a research attempt}
\author{GPT-6 (Codex)}
\date{9 October 2026}
\begin{document}
\maketitle
\section*{Target and outcome}
\textbf{Status: unresolved.} A local productivity gain does not determine employment, wages or task creation after equilibrium adjustment. This attempt proves that point in a competitive task-cost benchmark, derives a substitution-versus-scale condition with nonzero machine cost, and gives an explicit observational equivalence for the post-adoption invention path. It also defines a reproducible mechanism decomposition and worker-earnings target.

Brynjolfsson, Li and Raymond's primary working-paper record \cite{ai} concerns a workplace deployment and reports productivity heterogeneity. Its record distinguishes the 2023 working paper from the 2025 journal publication. Neither those local outcomes nor the numerical examples below are treated as an aggregate employment forecast.

\section*{Labor requirements, product demand and labor supply}
In one period, producing a unit of output requires $\ell>0$ paid human hours and machine inputs costing $k\geq0$ in an outside numeraire. Competitive constant-returns firms have unit cost and price
\[
 p=w\ell+k.
\]
Product demand is $Q=Kp^{-\varepsilon}$, with $K,\varepsilon>0$. Aggregate labor supply in hours is $L=Bw^\eta$, with $B,\eta>0$, and labor demand is $L=\ell Q$. Parameters and machine prices are fixed within the counterfactual. The product-demand schedule can represent an outside market, avoiding an unstated assumption that this one sector is the entire closed economy.

For any $\ell>0,k\geq0$, labor demand is continuous and strictly decreasing in $w>0$, while supply is strictly increasing. Demand exceeds supply for sufficiently small wages, and supply exceeds demand for sufficiently large wages. Thus there is a unique positive market-clearing wage. The equilibrium includes product-market clearing, zero profits at unit cost, and labor-market clearing.

When $k=0$, explicit solution gives
\[
 w=\left(\frac KB\ell^{1-\varepsilon}\right)^{1/(\eta+\varepsilon)},\qquad
 L=B\left(\frac KB\ell^{1-\varepsilon}\right)^{\eta/(\eta+\varepsilon)}.
\]
\begin{proposition}[A productivity improvement can raise or lower hours]
For a decline in $\ell$, paid hours rise if $\varepsilon>1$, fall if $\varepsilon<1$, and remain unchanged if $\varepsilon=1$ in the zero-machine-cost benchmark.
\end{proposition}
\begin{proof}
Log differentiation gives
$d\log L=\eta(1-\varepsilon)d\log\ell/(\eta+\varepsilon)$.
The denominator and $\eta$ are positive, while the productivity improvement has $d\log\ell<0$. The three cases follow.
\end{proof}

The fixed-output substitution effect is unambiguously negative: fewer hours are needed per unit. The equilibrium scale effect can dominate because prices fall and demand expands. Employment headcount is a separate variable; hours become proportional to headcount only after fixing hours per employed worker. This benchmark directly solves hours and wages, not all components of the catalogue trajectory.

\section*{An explicit identification counterexample}
Normalize $K=B=1$, initial $\ell=1$, $k=0$ and $\eta=1$. Economies with demand elasticities $\varepsilon=1/2$ and $\varepsilon=2$ both have initial $(w,p,Q,L)=(1,1,1,1)$. A workplace technology test at fixed output measures the same requirement change, $\ell:1\to0.8$, in both economies.

After economy-wide adoption, however,
\[
 L_{1/2}=0.8^{1/3}\approx0.9283,\qquad
 L_2=0.8^{-1/3}\approx1.0772.
\]
Thus the baseline equilibrium and local technical productivity effect are identical, but the equilibrium labor-hours signs differ. The models are observationally equivalent only for the stated baseline and workplace experiment; a suitable product-price experiment could identify their different demand curves. This specificity matters: the example does not claim equivalence for every possible dataset.

\section*{The scale condition with machine costs}
Let $s_L=w\ell/(w\ell+k)$ be the baseline labor share of unit cost. For local changes in $\ell,k$ with $k>0$, log differentiation of $L=\ell Kp^{-\varepsilon}$ and $L=Bw^\eta$ gives
\[
 (\eta+\varepsilon s_L)d\log w
 =(1-\varepsilon s_L)d\log\ell
   -\varepsilon(1-s_L)d\log k.
\]
Indeed, $d\log p=s_L(d\log w+d\log\ell)+(1-s_L)d\log k$, which substituted into demand and equated to supply proves the formula.

If machine cost is fixed, a reduction in human requirements raises hours exactly when $\varepsilon s_L>1$. If adoption lowers machine costs as well, the last term creates an additional expansion channel. A technology that merely changes task exposure without changing $\ell$, $k$ or feasible production has no effect in this model. Endogenous adoption costs, capital supply and entry would require further equations rather than an informal adjustment to this derivative.

\section*{Task creation is independently underidentified}
Let baseline adoption intensity be zero, and suppose new task labels are born at rate $\nu a_t$, where $\nu\geq0$ is unknown and labels come from a stable common taxonomy. Let active new-task stock obey
\[
 S_{t+1}=(1-\rho)S_t+\nu a_t,\qquad S_0=0,\quad 0<\rho\leq1.
\]
Before adoption, every value of $\nu$ produces exactly the same zero stock and zero new labels. Under constant $a_t=a>0$,
\[
 S_t=\nu a\frac{1-(1-\rho)^t}{\rho}.
\]
This follows by solving the first-order recurrence. Cumulative births through dates $0,\ldots,t-1$ equal $t\nu a$, a different object from the active stock because old tasks may disappear. Both can be measured only after fixing whether a relabeled task counts as new.

If compatible models allow $\nu\in[0,\bar\nu]$ independently of the baseline data, the displayed stock interval obtained by replacing $\nu$ by its endpoints is sharp: every intermediate value follows from a compatible $\nu$. Local productivity evidence on legacy tasks does not restrict it. Connecting these births to wages and hours requires a production and demand model for the new tasks; adding their count directly to jobs would be a units error.

\section*{Mechanism attribution and distributional earnings}
Define two technology components: a change in human requirements $d$ and an expansion or invention component $a$. For any equilibrium outcome $Y(d,a)$ define the interaction
\[
 I=Y(1,1)-Y(1,0)-Y(0,1)+Y(0,0).
\]
The total change is the sum of the two baseline marginal changes and $I$, by rearrangement. A symmetric attribution assigns half the interaction to each component; a sequential attribution assigns it to whichever component is applied second. These are declared conventions for a known four-counterfactual system. Observing only $Y(0,0)$ and $Y(1,1)$ does not identify the two missing cells or a unique mechanism allocation.

For predetermined worker type $x$, a meaningful earnings target over horizon $H$ is
\[
 \Delta V_x=\E\left[\sum_{t=0}^H\beta^t
 \{w_t(a)\,h_{it}(a)-w_t(a_0)\,h_{it}(a_0)\}\mid X_i=x\right],
\]
where paid hours $h_{it}$ equal zero during unemployment. Conditioning instead on being employed after adoption discards displacement and changes the population. Heterogeneous wages, skill transitions and household price effects are omitted from the scalar benchmark and must be supplied for this target.

The unresolved research task is to identify demand, adoption, invention and worker-transition primitives jointly, then solve the declared dynamic equilibrium under both regimes. The proved benchmark and counterexamples show exactly which conclusions cannot be obtained from task exposure or local productivity alone.

\section{Completion Estimate}
52\%

This is a post hoc, heuristic estimate for the full catalogue target.
A unique competitive equilibrium proves when productivity raises or lowers labor hours, and observationally equivalent elasticities and invention paths yield sharp restricted counterexamples. Dynamic adoption, entry, skill transitions, task-specific hours and baseline-worker earnings still require jointly identified primitives and equilibrium solutions.

\begin{thebibliography}{9}
\bibitem{ai} E. Brynjolfsson, D. Li and L. R. Raymond (2023), Generative AI at Work, NBER Working Paper 31161; journal publication 2025. Primary NBER abstract and version record inspected: \url{https://www.nber.org/papers/w31161}. Used for the scope of workplace evidence, not for an aggregate elasticity or task-birth estimate.
\end{thebibliography}
\end{document}
